\section{The critical number and cup length of the quadric}
\label{sec:critical-bound}

We now show that four is both the critical number and the rational cup
length of the ambient quadric. The lower bound is the classical
cup-product obstruction from Lusternik--Schnirelmann
theory~\cite{LusternikSchnirelmann1934}; see
Weber~\cite[Theorems~3.5 and 3.6]{Weber2019} for a modern treatment
valid without a Morse assumption. We give the differential-form proof,
then obtain the matching upper bound from the unequal-speed rotations
of Section~\ref{sec:quadric}.

\begin{lemma}\label{lem:critical-bound}
Let $M$ be a closed smooth manifold and $r\ge1$ an integer. Suppose there are classes
$a_j\in H^{d_j}(M;\R)$, with $d_j>0$ for $1\le j\le r$, such that
$a_1\smile\cdots\smile a_r\ne0$. Then every smooth function on $M$
has at least $r+1$ critical points.
\end{lemma}

\begin{proof}
Suppose that $h:M\to\R$ has at most $r$ critical points.
List its distinct critical values as $c_1<\cdots<c_k$, where $1\le k\le r$;
the absolute minimum and maximum are among them.
For a fixed Riemannian metric let $g_t$ be the negative-gradient flow
of $h$, which is complete by compactness, and put $P_b=\{h\le b\}$.
We first cover $M$ by $k$ open sets, each diffeomorphic to a finite
disjoint union of coordinate balls. Every closed positive-degree form
will be exact on each of these sets; the cover will therefore force
every product of $k$ positive-degree cohomology classes to vanish.

Around the critical points at $c_i$, choose an open set $U_i$ that is a
union of disjoint coordinate balls, and a smaller neighborhood $V_i$
with $\overline V_i\subset U_i$. A bound on the speed of the flow and
the positive distance between $M\setminus U_i$ and $\overline V_i$
give $T_i>0$ such that trajectories starting outside $U_i$ stay outside
$V_i$ for $0\le t\le T_i$.

Choose a closed level band $\{|h-c_i|\le\delta_i\}$ narrow enough that
it contains no critical points outside $V_i$. On the part outside
$V_i$, compactness gives a positive lower bound $b_i$ for
$|\nabla h|^2$; if that set is empty, any positive $b_i$ will do.
Choose $0<\epsilon_i<\delta_i$ so that the intervals
$[c_i-\epsilon_i,c_i+\epsilon_i]$ are separated and
$2\epsilon_i<T_i b_i$. Along the negative-gradient flow,
\[
 \frac{d}{dt}h(g_t(q))=-|\nabla h(g_t(q))|^2.
\]
A point of $P_{c_i+\epsilon_i}\setminus U_i$ either already lies below
$c_i-\epsilon_i$ or decreases at rate at least $b_i$ until it reaches
that level. It follows that
\begin{equation}\label{eq:sublevel-drop}
 g_{T_i}(P_{c_i+\epsilon_i}\setminus U_i)
       \subset P_{c_i-\epsilon_i}.
\end{equation}
For $i>1$, a further fixed flow time $S_i>0$ carries $P_{c_i-\epsilon_i}$
into $P_{c_{i-1}+\epsilon_{i-1}}$: the compact band between these levels
contains no critical points and has a positive gradient-norm lower bound.

Since $c_1$ is the absolute minimum, $P_{c_1-\epsilon_1}$ is empty,
and \eqref{eq:sublevel-drop} shows that $P_{c_1+\epsilon_1}\subset U_1$.
Inductively, cover $P_{c_i+\epsilon_i}$ by $U_i$ and the preimages under
$g_{T_i+S_i}$ of the $i-1$ sets covering the preceding upper sublevel.
Equation~\eqref{eq:sublevel-drop} proves that these sets cover.
Each $g_t$ is a diffeomorphism of $M$, so every member remains
diffeomorphic to a finite disjoint union of coordinate balls. The last
upper sublevel is all of $M$, since $c_k$ is the absolute maximum.

Write the resulting cover as $O_1,\ldots,O_k$. Choose closed forms
$\alpha_j$ representing $a_j$ for $1\le j\le k$. Since these forms
have positive degree, the Poincar\'e lemma on each ball gives
$\alpha_j=d\lambda_j$ on $O_j$. Choose a smooth partition of unity
$\psi_1,\ldots,\psi_k$ subordinate to the cover. Let
$\chi:[0,1]\to[0,1]$ be smooth, zero on $[0,1/(4k)]$ and one on
$[1/(2k),1]$, and put $\rho_j=\chi(\psi_j)$. At every point some
$\psi_j\ge1/k$, so some $\rho_j$ is identically one nearby.
Each $\rho_j\lambda_j$ extends smoothly by zero outside $O_j$.

The closed forms
\[
 \alpha'_j=\alpha_j-d(\rho_j\lambda_j)
\]
represent the same classes as $\alpha_j$. Near every point, at least one
of these forms vanishes, and hence $\alpha'_1\wedge\cdots\wedge\alpha'_k=0$.
Thus $a_1\smile\cdots\smile a_k=0$, contradicting the nonvanishing
of $a_1\smile\cdots\smile a_r$ because $k\le r$.
\end{proof}

Extension of coefficients $H^*(M;\mathbb Q)\to H^*(M;\R)$ is injective
for a closed manifold and respects cup products. The lemma therefore
also gives
\begin{equation}\label{eq:cup-critical}
 \operatorname{cuplength}(M;\mathbb Q)\le\Crit(M),
\end{equation}
with the unit-inclusive convention of Section~\ref{sec:introduction}.
It also recovers the symplectic-volume obstruction used in our example.

\begin{corollary}\label{cor:symplectic-critical}
Every smooth function on a closed connected symplectic manifold of real
dimension $2n$, with $n\ge1$, has at least $n+1$ critical points.
\end{corollary}

\begin{proof}
The class $[\omega]^n$ is nonzero, since
$\int_M\omega^n>0$ in the symplectic orientation. Apply
Lemma~\ref{lem:critical-bound} with each $a_j=[\omega]$.
\end{proof}

\begin{corollary}\label{cor:quadric-critical}
For the quadric threefold in \eqref{eq:quadric},
\[
 \Crit(Q^3)=\operatorname{cuplength}(Q^3;\mathbb Q)=4.
\]
\end{corollary}

\begin{proof}
Let $h=c_1(\mathcal O_{Q^3}(1))\in H^2(Q^3;\mathbb Z)$ be the
hyperplane class. Since $Q^3$ has degree two in $\CP^4$,
\[
 h^3=2[\mathrm{pt}],
\]
where $[\mathrm{pt}]$ evaluates to one on the fundamental class.
This product remains nonzero over $\mathbb Q$, giving
$\operatorname{cuplength}(Q^3;\mathbb Q)\ge4$.

For the matching upper bound, take the Hamiltonian $G_{1,2}$ of
\eqref{eq:rotation-hamiltonian}. Its critical points are precisely the
zeros of its Hamiltonian vector field. That field is induced on $Q^3$
by the complex-linear generator
\[
 T=\operatorname{diag}(0,J,2J),
 \qquad J=\begin{pmatrix}0&-1\\1&0\end{pmatrix}.
\]
The induced projective vector field vanishes at $[z]$ exactly when
$Tz\in\C z$. The eigenvalues $0,\pm i,\pm2i$ of $T$ are distinct,
so its projective zeros are its five eigenlines. The zero eigenline
$[1:0:0:0:0]$ misses $Q^3$, whereas the other four lines
\[
 [0:1:\pm i:0:0],\qquad [0:0:0:1:\pm i]
\]
all satisfy the quadric equation. Thus $G_{1,2}$ has exactly four
critical points on $Q^3$. Together with \eqref{eq:cup-critical}, this gives
\[
 4\le\operatorname{cuplength}(Q^3;\mathbb Q)
   \le\Crit(Q^3)\le4.
\]
\end{proof}
